\chapter{Despliegue}\label{cap:despliegue}

\section{Docker Compose}

El fichero \path{docker-compose.yml} define dos servicios:

\begin{itemize}[leftmargin=*,itemsep=0.2em]
  \item \texttt{backend}: imagen construida desde
        \path{fakenews-backend/Dockerfile}, expone el puerto 8000.
  \item \texttt{frontend}: imagen construida desde
        \path{fakenews-frontend/Dockerfile} con multi-stage build
        (Vite \emph{build} + \texttt{nginx}), expone el puerto 5173
        (mapeado a 80 / 443 en producción).
\end{itemize}

Las variables de entorno se cargan desde \path{.env} (no
versionado).

\section{Imágenes}

\paragraph{Backend.} Imagen base \texttt{python:3.12-slim};
instalación de dependencias en una capa cacheada antes del código,
para minimizar invalidación de caché en cada cambio fuente.

\paragraph{Frontend.} Multi-stage:
\begin{enumerate}[leftmargin=*,itemsep=0.2em]
  \item \texttt{node:20-alpine} ejecuta \texttt{npm ci} y
        \texttt{npm run build}.
  \item \texttt{nginx:alpine} sirve los artefactos estáticos con la
        configuración de \path{nginx.conf} (rewrites para SPA,
        cabeceras de seguridad).
\end{enumerate}

\section{Variables de entorno}

Documentadas en \path{.env.example}. Las principales:
\texttt{SUPABASE\_URL}, \texttt{SUPABASE\_ANON\_KEY},
\texttt{SUPABASE\_SERVICE\_ROLE\_KEY},
\texttt{STRIPE\_SECRET\_KEY}, \texttt{STRIPE\_WEBHOOK\_SECRET},
\texttt{TAVILY\_API\_KEY},
\texttt{ALLOWED\_ORIGINS}.

\section{Despliegue en VPS}

\begin{enumerate}[leftmargin=*,itemsep=0.2em]
  \item Configurar dominio y certificado TLS (Let's Encrypt vía
        \texttt{certbot} o reverse-proxy Caddy).
  \item Pull del repositorio y \texttt{docker compose up -d --build}.
  \item Configurar \emph{webhook} de Stripe apuntando a
        \texttt{https://api.dominio/billing/webhook}.
  \item Aplicar migraciones SQL desde la consola Supabase.
\end{enumerate}

\section{Observabilidad}

\begin{itemize}[leftmargin=*,itemsep=0.2em]
  \item \emph{Logs} estructurados a \texttt{stdout}; en producción se
        recogen con la pila estándar de Docker (\texttt{docker logs}).
  \item \texttt{GET /health} para sondas de \emph{liveness}.
  \item Métricas de uso almacenadas en la tabla \texttt{usage} y
        consultables desde el panel administrativo.
\end{itemize}
